\documentclass[11pt]{article}
\usepackage{amsmath,amssymb}
\title{The retained conditional residual as a Gaussian score contact}\author{}\date{}
\begin{document}\maketitle
\paragraph{Frame and input.}
Use the same-coefficient conditioning identity of the same coefficient schur cumulant decomposition calculation and its
harmonic boundary response. Arbitrary internal insertions here are
DEFINED in that common-tree fixed-log frame. Equality of the scalar
coefficient to a local FP representation does not identify arbitrary
internal source operators across nonlinear coordinate changes.
Such an identification must carry its transformed source and Jacobian
partners; it is not supplied by this note.

Fix one complete block, its actual free tree coordinates $\eta$, and
its original conditional flat covariance $C=A^{-1}$. Boundary action
indices are retained, while only genuinely integrated links carry Haar.
Define the FULL local polynomial
\[
 P(z)=T_4(z)-J_{2,I}(z)-\tfrac12T_3(z)^2.                       \tag{1}
\]
The gauge/ghost effects of the original tree integral are implicit in
this exact choice of coordinates, not discarded from a different
representation. At flat SU(2), the color-diagonal covariance gives
$ET_3=0$ and $ED_wT_3=0$ for every deterministic direction $w$,
including boundary directions. Consequently the local variance-defined
second insertion tensor is exactly
\[
 K_C(u,v)=E_C D_uD_v P(\eta).                                 \tag{2}
\]
This equality would require mean partners when the tadpole is nonzero.

\paragraph{1. Exact cross and residual contacts.}
Let $h$ be the true local harmonic extension of a retained boundary
direction and $r$ a zero-boundary residual in this same frame. Put
$q=Ar$, $e_r=r^*Ar$, and $\ell_r=q^*\eta$. Then Gaussian integration
by parts, with the boundary coordinates held fixed, gives
\[
 \boxed{K_C(h,r)=E_C[\ell_r D_hP],\qquad
 K_C(r,r)=E_C[(\ell_r^2-e_r)P].}                              \tag{3}
\]
No tadpole, Haar, quartic or cubic cross term of (1) is removed from
these expectations. The second identity has the subtraction $-e_r$;
retaining only $\ell_r^2$ would introduce a spurious vacuum contact.
Both identities are exact even if $h$ has interior components: the
constant differential operators commute, and integration by parts is
only in the actual zero-boundary direction $r$.

Since $E\ell_r=0$ and $E\ell_r^2=e_r$, (3) can equivalently be written
as centered covariances. In particular,
\[
 |K_C(h,r)|\le\sqrt{e_r}\,
                \|D_hP-E D_hP\|_2,
\]
\[
 |K_C(r,r)|\le\sqrt2\,e_r\,\|P-EP\|_2.                       \tag{4}
\]
These are exact quantitative interfaces, not claims that the displayed
polynomial norms are uniformly bounded. They retain the complete
connected polynomial. Estimating its separate vacuum summands by
absolute values can lose the desired scale behavior.

\paragraph{2. Relation to the actual long-leg gain.}
For a deterministic marked field $w$,54 constructs the true local
residual with
\[
 \sum_B e_{r_B}\le \ell^2\|Hw\|^2/12.
\]
A linear relative gauge transformation to the declared common-tree
coordinates preserves these classical energies. Thus the score
variance in (3) is EXACTLY the residual energy controlled by the long leg harmonic residual with gauss constraint calculation.
No tree Euclidean norm is substituted for the full-link angular norm.
After the local polynomial average is taken, one may average the
marked direction against an independent heat-covariance realization;
54 then supplies the extra $\ell^2/(1+s)$ factor for its energy.
One must not use (3) by treating a direction correlated with $\eta$
as deterministic.

This identifies the missing response norm precisely: the complete
conditional fluctuations in (4), or a signed estimate combining their
expectations with the induced coarse terms of the same coefficient schur cumulant decomposition calculation. The residual-energy gain
alone does not bound those fluctuations. Nor does harmonicity set
$K_C(h,r)$ to zero: (3) displays its surviving score insertion.

\paragraph{3. Same-coefficient compensation still remains.}
Under the original joint skeleton law,85 retains
\[
 T_4^{\rm eff}=T_4(Hy)
 -\tfrac12\nabla_I T_3(Hy)^*C_I\nabla_I T_3(Hy)
\]
and the FULL $\operatorname{Var}_{S^{-1}}T_3(Hy)$, including all
cross-block contractions. These are the compensation for replacing a
global covariance by local conditional covariances; they must accompany
(3). No equality between the local sum and the global self-energy is
asserted without that remainder. Its source derivatives and any change
to the full vector/ghost frame remain explicit work.

The identities isolate genuine original conditional score contacts,
but do not establish a logarithmic bound for the complete first
nonlinear coefficient or its original uniform remainder. No physical
scale iteration or S94 claim follows.
\end{document}
